\chapter{Del gen a la proteína}\label{ch:g11:gene-expression}

En 1961, un bioquímico alimentó un extracto acelular de bacterias con
un \omterm{def:g11:gene-expression:transcription}{ARN} artificial hecho solo con la letra U, repetida. El extracto
fabricó una \omterm{def:g11:gene-expression:protein}{proteína} hecha solo con un \omterm{def:g10:chemistry-of-life:families}{aminoácido}, la fenilalanina,
repetido. Con ese único experimento se leyó la primera palabra del
\omterm{def:g11:gene-expression:code}{código genético}: UUU significa fenilalanina. En cinco años se conocían
las sesenta y cuatro palabras, y resultaron ser las mismas en una
bacteria, en una planta de trigo y en un humano. Este capítulo sigue el
camino que va de la secuencia de \omterm{def:g10:universal-dna:nucleotide}{nucleótidos} de un \omterm{def:g10:universal-dna:gene}{gen} a la secuencia
de \omterm{def:g10:chemistry-of-life:families}{aminoácidos} de una \omterm{def:g11:gene-expression:protein}{proteína} --- el camino que recorre todo rasgo del
\cref{ch:g11:mutations}.

\section{Los genes fabrican proteínas}

\begin{proposition}[Un gen, una proteína]\label{prop:g11:gene-expression:onegene}
Un \omterm{def:g10:universal-dna:gene}{gen} se expresa cuando la \omterm{def:g10:cells-common-unit:cell}{célula} usa su secuencia para fabricar la
\omterm{def:g11:gene-expression:protein}{proteína} que codifica. Casi todos los rasgos dependen de \omterm{def:g10:chemistry-of-life:families}{proteínas} ---
\omterm{def:g10:cell-metabolism:metabolism}{enzimas}, fibras estructurales, transportadores, receptores --- y una
\omterm{def:g11:mutations:mutation}{mutación} cambia un rasgo cambiando la \omterm{def:g11:gene-expression:protein}{proteína} que produce su \omterm{def:g10:universal-dna:gene}{gen}.
\end{proposition}
\begin{proof}[Pruebas experimentales]
Beadle y Tatum (1941) irradiaron esporas de un moho del pan y
recogieron mutantes que ya no podían crecer en medio mínimo a menos que
se añadiera una sustancia concreta, como el \omterm{def:g10:chemistry-of-life:families}{aminoácido} arginina. A cada
mutante le faltaba una \omterm{def:g10:cell-metabolism:metabolism}{enzima} de la cadena de reacciones que fabrica
esa sustancia; cada defecto se heredaba como un solo \omterm{def:g10:universal-dna:gene}{gen}; y los
mutantes bloqueados en pasos distintos de la misma cadena llevaban
\omterm{def:g11:mutations:mutation}{mutaciones} en \omterm{def:g10:universal-dna:gene}{genes} distintos. Un \omterm{def:g10:universal-dna:gene}{gen}, una \omterm{def:g10:cell-metabolism:metabolism}{enzima} --- y, como
generalizaron después otros trabajos, un \omterm{def:g10:universal-dna:gene}{gen}, una cadena polipeptídica.
\end{proof}

\begin{definition}[Proteína, secuencia de aminoácidos]\label{def:g11:gene-expression:protein}
Una \emph{proteína}\index{proteína} es una cadena de \omterm{def:g10:chemistry-of-life:families}{aminoácidos}, de
unas pocas decenas a varios miles de unidades, tomados de un conjunto
de veinte clases. La \emph{secuencia} --- el orden de los \omterm{def:g10:chemistry-of-life:families}{aminoácidos}
--- es lo que especifica el \omterm{def:g10:universal-dna:gene}{gen}; una vez fabricada, la cadena se pliega
en una forma tridimensional definida que determina esa secuencia, y la
forma determina la función. Cambia un \omterm{def:g10:chemistry-of-life:families}{aminoácido} y pueden cambiar el
plegamiento y la función.
\end{definition}

\section{La transcripción: el gen copiado en ARN}

\begin{definition}[ARN mensajero y transcripción]\label{def:g11:gene-expression:transcription}
El \emph{ARN}\index{ARN} (ácido ribonucleico) es una cadena de una sola
hebra de \omterm{def:g10:universal-dna:nucleotide}{nucleótidos} como los del \omterm{def:g10:universal-dna:information}{ADN}, salvo que su azúcar es la ribosa
y la base uracilo (U) sustituye a la timina (T). La
\emph{transcripción}\index{transcripción} es la copia de la secuencia
de un \omterm{def:g10:universal-dna:gene}{gen} en una molécula de \emph{ARN mensajero}\index{ARN mensajero}
(ARNm): la \omterm{def:g10:cell-metabolism:metabolism}{enzima} \emph{ARN polimerasa} abre la \omterm{prop:g10:universal-dna:helix}{doble hélice} a lo largo
del \omterm{def:g10:universal-dna:gene}{gen}, lee una hebra (la \emph{hebra molde}) y ensambla el ARN
complementario, con U frente a A, A frente a T, G frente a C y C frente
a G. El ARNm tiene la secuencia de la otra hebra del \omterm{def:g10:universal-dna:gene}{gen}, con U en
lugar de T; sale del \omterm{def:g10:cells-common-unit:organelle}{núcleo} hacia el \omterm{def:g10:cells-common-unit:cell}{citoplasma}.
\end{definition}

\begin{omfigure}
\begin{tikzpicture}[font=\small, >=Stealth, line cap=round]
  % DNA strands with bubble
  \draw[very thick, omDef] (0,0.5) -- (2.5,0.5) .. controls (3.2,0.5) and (3.4,1.3) .. (4.3,1.3) -- (6.7,1.3) .. controls (7.6,1.3) and (7.8,0.5) .. (8.5,0.5) -- (11,0.5);
  \draw[very thick, omDef] (0,-0.5) -- (2.5,-0.5) .. controls (3.2,-0.5) and (3.4,-1.3) .. (4.3,-1.3) -- (6.7,-1.3) .. controls (7.6,-1.3) and (7.8,-0.5) .. (8.5,-0.5) -- (11,-0.5);
  \foreach \x in {0.3,0.7,...,2.3} \draw[gray] (\x,0.5) -- (\x,-0.5);
  \foreach \x in {8.7,9.1,...,10.9} \draw[gray] (\x,0.5) -- (\x,-0.5);
  % template strand letters (bottom strand in bubble)
  \foreach \x/\l in {4.5/T, 4.9/A, 5.3/C, 5.7/G, 6.1/G, 6.5/A} \node[font=\scriptsize, omDef, anchor=north] at (\x,-1.32) {\l};
  % mRNA
  \draw[very thick, omThm] (4.4,-0.9) -- (6.6,-0.9) .. controls (7.4,-0.9) and (7.2,0.1) .. (8.0,0.6) .. controls (8.6,1.0) and (9.4,1.4) .. (10.6,1.7);
  \foreach \x/\l in {4.5/A, 4.9/U, 5.3/G, 5.7/C, 6.1/C, 6.5/U} \node[font=\scriptsize, omThm, anchor=south] at (\x,-0.85) {\l};
  \foreach \x in {4.5,4.9,...,6.5} \draw[gray, densely dotted] (\x,-0.95) -- (\x,-1.28);
  % polymerase
  \draw[thick, fill=omProp!30, rounded corners=8pt] (6.95,-1.6) rectangle (8.65,0.2);
  \node[align=center] at (7.8,-0.7) {ARN\\ polimerasa};
  \draw[->, thick] (8.7,-1.9) -- (9.8,-1.9) node[right] {avanza por el gen};
  \node[anchor=east, omDef, align=right] at (-0.15,0) {ADN del\\ gen};
  \node[anchor=west, omThm, align=left] at (10.7,1.7) {ARNm, que se\\ va liberando};
  \node[anchor=north, omDef] at (5.5,-1.75) {hebra molde};
\end{tikzpicture}

{\small La \omterm{def:g11:gene-expression:transcription}{transcripción}. La \omterm{def:g11:gene-expression:transcription}{ARN} polimerasa abre la hélice, lee la
hebra molde y construye un \omterm{def:g11:gene-expression:transcription}{ARN mensajero} complementario de ella (con U
en lugar de T). Detrás de la \omterm{def:g10:cell-metabolism:metabolism}{enzima} la hélice se vuelve a cerrar; el
\omterm{def:g11:gene-expression:transcription}{ARN} se despega.}
\end{omfigure}

\begin{example}[Un transcrito]\label{ex:g11:gene-expression:transcript}
La hebra molde \texttt{3'-TAC GGA CTT ATC-5'} da el ARNm
\texttt{5'-AUG CCU GAA UAG-3'}. La otra hebra del \omterm{def:g10:universal-dna:gene}{gen} dice
\texttt{5'-ATG CCT GAA TAG-3'}: el ARNm es la secuencia de esa hebra
con U en lugar de T, y por eso las secuencias de los \omterm{def:g10:universal-dna:gene}{genes} se escriben
por convenio como esa hebra, la \emph{hebra codificante}.
\end{example}

\section{El código genético}

\begin{definition}[Codón y código genético]\label{def:g11:gene-expression:code}
El ARNm se lee en grupos consecutivos de tres \omterm{def:g10:universal-dna:nucleotide}{nucleótidos}, los
\emph{codones}\index{codón}, a partir de un punto de partida definido.
El \emph{código genético}\index{código genético} es la correspondencia
entre los 64 codones posibles y los 20 \omterm{def:g10:chemistry-of-life:families}{aminoácidos}: 61 codones
especifican un \omterm{def:g10:chemistry-of-life:families}{aminoácido}, tres (UAA, UAG, UGA) son señales de
\emph{stop} que terminan la cadena, y AUG, que especifica la metionina,
sirve además de codón de \emph{inicio}. El código es
\emph{degenerado} --- casi todos los \omterm{def:g10:chemistry-of-life:families}{aminoácidos} tienen varios codones
--- y \emph{universal}: la misma tabla sirve para todos los organismos
conocidos, con raras variantes menores.
\end{definition}

\begin{omfigure}
\begin{tabular}{c|cccc|c}
\multicolumn{1}{c}{} & \multicolumn{4}{c}{\small segunda letra} & \multicolumn{1}{c}{} \\
{\small primera} & U & C & A & G & {\small tercera} \\ \hline
U & UUU Phe & UCU Ser & UAU Tyr & UGU Cys & U \\
 & UUC Phe & UCC Ser & UAC Tyr & UGC Cys & C \\
 & UUA Leu & UCA Ser & UAA \textbf{stop} & UGA \textbf{stop} & A \\
 & UUG Leu & UCG Ser & UAG \textbf{stop} & UGG Trp & G \\ \hline
C & CUU Leu & CCU Pro & CAU His & CGU Arg & U \\
 & CUC Leu & CCC Pro & CAC His & CGC Arg & C \\
 & CUA Leu & CCA Pro & CAA Gln & CGA Arg & A \\
 & CUG Leu & CCG Pro & CAG Gln & CGG Arg & G \\ \hline
A & AUU Ile & ACU Thr & AAU Asn & AGU Ser & U \\
 & AUC Ile & ACC Thr & AAC Asn & AGC Ser & C \\
 & AUA Ile & ACA Thr & AAA Lys & AGA Arg & A \\
 & AUG \textbf{Met} & ACG Thr & AAG Lys & AGG Arg & G \\ \hline
G & GUU Val & GCU Ala & GAU Asp & GGU Gly & U \\
 & GUC Val & GCC Ala & GAC Asp & GGC Gly & C \\
 & GUA Val & GCA Ala & GAA Glu & GGA Gly & A \\
 & GUG Val & GCG Ala & GAG Glu & GGG Gly & G \\
\end{tabular}

{\small El \omterm{def:g11:gene-expression:code}{código genético}, leído sobre el ARNm. Cada \omterm{def:g11:gene-expression:code}{codón} se
encuentra por su primera letra (fila), su segunda letra (columna) y su
tercera letra (línea dentro del bloque). AUG es a la vez metionina y la
señal de inicio; tres \omterm{def:g11:gene-expression:code}{codones} son de stop.}
\end{omfigure}

\begin{proposition}[Cómo se leyó el código]\label{prop:g11:gene-expression:readcode}
El significado de cada \omterm{def:g11:gene-expression:code}{codón} se estableció experimentalmente.
\end{proposition}
\begin{proof}[Pruebas experimentales]
Nirenberg y Matthaei (1961) añadieron \omterm{def:g11:gene-expression:transcription}{ARN} sintético de un solo
\omterm{def:g10:universal-dna:nucleotide}{nucleótido} repetido a un extracto acelular con \omterm{def:g10:cells-common-unit:organelle}{ribosomas}, \omterm{def:g10:chemistry-of-life:families}{aminoácidos} y
las \omterm{def:g10:cell-metabolism:metabolism}{enzimas} de la síntesis de \omterm{def:g10:chemistry-of-life:families}{proteínas}: el poli-U dio una cadena de
fenilalaninas, el poli-C una cadena de prolinas y el poli-A de lisinas.
Los \omterm{def:g11:gene-expression:transcription}{ARN} sintéticos de parejas y de tripletes repetidos y, después, la
unión de tripletes sueltos a los \omterm{def:g10:cells-common-unit:organelle}{ribosomas} asignaron los \omterm{def:g11:gene-expression:code}{codones}
restantes hacia 1966. La longitud de tres letras se había deducido de
la aritmética (dos letras dan solo 16 combinaciones, muy pocas para 20
\omterm{def:g10:chemistry-of-life:families}{aminoácidos}; tres dan 64) y se confirmó con \omterm{def:g11:mutations:mutation}{mutaciones}: insertar uno o
dos \omterm{def:g10:universal-dna:nucleotide}{nucleótidos} en un \omterm{def:g10:universal-dna:gene}{gen} destruye su \omterm{def:g11:gene-expression:protein}{proteína}, e insertar tres
restituye una casi normal.
\end{proof}

\begin{method}[Traducir una secuencia]\label{met:g11:gene-expression:translate}
\begin{enumerate}
  \item Escribe la hebra codificante (o transcribe la hebra molde):
        sustituye T por U para obtener el ARNm.
  \item Busca el primer AUG: es el inicio, y a partir de él queda fijado
        el marco de lectura.
  \item Corta el ARNm en tripletes consecutivos desde el AUG y lee cada
        uno en la tabla.
  \item Detente en el primer \omterm{def:g11:gene-expression:code}{codón} de stop; los \omterm{def:g10:chemistry-of-life:families}{aminoácidos} anotados,
        en orden, son la secuencia de la \omterm{def:g11:gene-expression:protein}{proteína}.
  \item Para un mutante, repite y compara: misma \omterm{def:g11:gene-expression:protein}{proteína}
        (silenciosa), un \omterm{def:g10:chemistry-of-life:families}{aminoácido} cambiado (de cambio de sentido),
        stop prematuro (sin sentido), todo cambiado a partir de la
        \omterm{def:g11:mutations:mutation}{mutación} (desplazamiento del marco de lectura).
\end{enumerate}
\end{method}

\begin{example}[Cuatro mutaciones, cuatro resultados]\label{ex:g11:gene-expression:outcomes}
Hebra codificante \texttt{ATG CCT GAA TTC TAG}: ARNm
\texttt{AUG CCU GAA UUC UAG}, \omterm{def:g11:gene-expression:protein}{proteína} Met--Pro--Glu--Phe.
\begin{itemize}
  \item \texttt{ATG CC\textbf{C} GAA TTC TAG}: CCC sigue siendo Pro ---
        \emph{silenciosa}.
  \item \texttt{ATG CCT G\textbf{C}A TTC TAG}: GCA es Ala ---
        Met--Pro--Ala--Phe, \emph{de cambio de sentido}.
  \item \texttt{ATG CCT \textbf{T}AA TTC TAG}: UAA es stop ---
        Met--Pro, \emph{sin sentido}, una \omterm{def:g11:gene-expression:protein}{proteína} truncada.
  \item \texttt{ATG CC\textbf{A}T GAA TTC TAG}: una A insertada
        desplaza el marco: \texttt{AUG CCA UGA}\dots{} Met--Pro--stop
        --- \emph{desplazamiento del marco}.
\end{itemize}
\end{example}

\section{La traducción: el mensaje leído}

\begin{proposition}[La traducción]\label{prop:g11:gene-expression:translation}
La \emph{traducción}\index{traducción} tiene lugar en los
\emph{\omterm{def:g10:cells-common-unit:organelle}{ribosomas}}, en el \omterm{def:g10:cells-common-unit:cell}{citoplasma}. Un \omterm{def:g10:cells-common-unit:organelle}{ribosoma} se une al ARNm por su
\omterm{def:g11:gene-expression:code}{codón} de inicio y avanza por él \omterm{def:g11:gene-expression:code}{codón} a \omterm{def:g11:gene-expression:code}{codón}. A cada \omterm{def:g11:gene-expression:code}{codón} le
corresponde un \emph{\omterm{def:g11:gene-expression:transcription}{ARN} de transferencia} (ARNt), un \omterm{def:g11:gene-expression:transcription}{ARN} pequeño que
lleva, en un extremo, el \emph{anticodón} de tres \omterm{def:g10:universal-dna:nucleotide}{nucleótidos}
complementario del \omterm{def:g11:gene-expression:code}{codón} y, en el otro, el \omterm{def:g10:chemistry-of-life:families}{aminoácido} correspondiente.
El \omterm{def:g10:cells-common-unit:organelle}{ribosoma} une cada \omterm{def:g10:chemistry-of-life:families}{aminoácido} que llega a la cadena en crecimiento y
suelta el ARNt vacío; en un \omterm{def:g11:gene-expression:code}{codón} de stop libera la cadena terminada,
que se pliega. Varios \omterm{def:g10:cells-common-unit:organelle}{ribosomas} leen un mismo ARNm uno tras otro, y un
solo ARNm rinde cientos de copias de la \omterm{def:g11:gene-expression:protein}{proteína} antes de degradarse.
\end{proposition}
\admitted

\begin{omfigure}
\begin{tikzpicture}[font=\small, >=Stealth, line cap=round]
  % mRNA
  \draw[very thick, omThm] (0,0) -- (11,0);
  \foreach \x/\l in {0.4/A,0.8/U,1.2/G, 1.7/C,2.1/C,2.5/U, 3.0/G,3.4/A,3.8/A, 4.3/U,4.7/U,5.1/C, 5.6/G,6.0/G,6.4/U, 6.9/U,7.3/A,7.7/G}
    \node[font=\scriptsize, omThm, anchor=north] at (\x,-0.05) {\l};
  \node[anchor=east, omThm] at (-0.15,0) {ARNm};
  % ribosome (two lobes) over codons 3 and 4
  \draw[thick, fill=omProp!25, rounded corners=10pt] (2.7,0.15) rectangle (5.5,1.4);
  \draw[thick, fill=omProp!35, rounded corners=12pt] (2.5,1.4) .. controls (2.5,3.0) and (5.7,3.0) .. (5.7,1.4) -- cycle;
  \node at (4.1,2.1) {ribosoma};
  % tRNAs
  \draw[very thick, omMeth!80!black] (3.4,0.55) -- (3.4,1.35);
  \node[font=\scriptsize, omMeth!80!black, anchor=south] at (3.4,0.15) {CUU};
  \draw[very thick, omMeth!80!black] (4.7,0.55) -- (4.7,1.35);
  \node[font=\scriptsize, omMeth!80!black, anchor=south] at (4.7,0.15) {AAG};
  % amino acids chain
  \foreach \i/\x in {0/3.4, 1/4.7} \fill[omDef!70] (\x,1.5) circle (4pt);
  \foreach \x in {2.6, 1.9, 1.2} \fill[omDef!70] (\x,1.9) circle (4pt);
  \draw[thick, omDef!70] (1.2,1.9) -- (2.6,1.9) -- (3.4,1.5) -- (4.7,1.5);
  \node[anchor=east, omDef!80!black, align=right] at (0.9,1.9) {proteína en\\ crecimiento};
  % incoming tRNA
  \draw[very thick, omMeth!80!black] (7.0,1.3) -- (7.0,2.1);
  \fill[omDef!70] (7.0,2.25) circle (4pt);
  \node[font=\scriptsize, omMeth!80!black, anchor=west] at (7.15,1.35) {CCA};
  \node[anchor=west, align=left, font=\scriptsize] at (7.3,1.7) {llega el siguiente ARNt\\ con su aminoácido};
  \draw[->, thick] (7.0,1.1) -- (6.1,0.75);
  \draw[->, thick] (5.7,-0.9) -- (7.2,-0.9) node[right] {el ribosoma avanza};
  \node[anchor=north, font=\scriptsize] at (4.05,-0.45) {codón};
  \draw (2.85,-0.35) -- (2.85,-0.45) -- (3.95,-0.45) -- (3.95,-0.35);
\end{tikzpicture}

{\small La \omterm{prop:g11:gene-expression:translation}{traducción}. El \omterm{def:g10:cells-common-unit:organelle}{ribosoma} abarca dos \omterm{def:g11:gene-expression:code}{codones}; un ARNt cuyo
anticodón encaja con cada \omterm{def:g11:gene-expression:code}{codón} trae su \omterm{def:g10:chemistry-of-life:families}{aminoácido}, la cadena se
transfiere al recién llegado y el \omterm{def:g10:cells-common-unit:organelle}{ribosoma} avanza un \omterm{def:g11:gene-expression:code}{codón}. En un \omterm{def:g11:gene-expression:code}{codón}
de stop la cadena se libera.}
\end{omfigure}

\begin{example}[Velocidad y rendimiento]\label{ex:g11:gene-expression:speed}
Un \omterm{def:g10:cells-common-unit:organelle}{ribosoma} añade de 5 a 20 \omterm{def:g10:chemistry-of-life:families}{aminoácidos} por segundo: una \omterm{def:g11:gene-expression:protein}{proteína} de
300 \omterm{def:g10:chemistry-of-life:families}{aminoácidos} tarda menos de un minuto. Los \omterm{def:g10:cells-common-unit:organelle}{ribosomas} se siguen unos
a otros sobre un ARNm separados unos 80 \omterm{def:g10:universal-dna:nucleotide}{nucleótidos}, así que un mensaje
de 900 \omterm{def:g10:universal-dna:nucleotide}{nucleótidos} lleva una docena a la vez; a lo largo de su vida de
unas horas, un solo ARNm produce mil moléculas de \omterm{def:g11:gene-expression:protein}{proteína}. Una \omterm{def:g10:cells-common-unit:cell}{célula}
de \emph{E.~coli} fabrica unas $\num{20000}$ moléculas de \omterm{def:g11:gene-expression:protein}{proteína} por
segundo; el precursor de un glóbulo rojo, unas $\num{5e8}$ moléculas de
hemoglobina en toda su vida.
\end{example}

\section{Un gen, varios mensajes}

\begin{proposition}[Maduración del mensaje en los eucariotas]\label{prop:g11:gene-expression:splicing}
En las \omterm{def:g10:cells-common-unit:prokaryote}{células eucariotas}, la secuencia de un \omterm{def:g10:universal-dna:gene}{gen} está interrumpida por
\emph{intrones}\index{intrón}, segmentos que se transcriben pero que
después se recortan del \omterm{def:g11:gene-expression:transcription}{ARN} antes de que salga del \omterm{def:g10:cells-common-unit:organelle}{núcleo}; los
segmentos que se conservan y se unen extremo con extremo, los
\emph{exones}, forman el ARNm maduro. Muchos \omterm{def:g10:universal-dna:gene}{genes} pueden recortarse de
más de una manera y conservar juegos distintos de exones
(\emph{maduración alternativa}), de modo que un solo \omterm{def:g10:universal-dna:gene}{gen} rinde varios
ARNm distintos y varias \omterm{def:g10:chemistry-of-life:families}{proteínas} emparentadas: los $\num{20000}$ \omterm{def:g10:universal-dna:gene}{genes}
humanos codifican bastante más de $\num{100000}$ \omterm{def:g10:chemistry-of-life:families}{proteínas}.
\end{proposition}
\admitted

\begin{omfigure}
\begin{tikzpicture}[font=\small, >=Stealth]
  \tikzset{ex/.style={draw, fill=omDef!40, minimum height=0.45cm, minimum width=1.0cm, inner sep=1pt},
           intr/.style={draw, fill=gray!20, minimum height=0.45cm, minimum width=0.7cm, inner sep=1pt}}
  \node[anchor=east] at (-0.2,0) {gen (ADN)};
  \node[ex] at (0.5,0) {1}; \node[intr] at (1.35,0) {}; \node[ex] at (2.2,0) {2}; \node[intr] at (3.05,0) {}; \node[ex] at (3.9,0) {3}; \node[intr] at (4.75,0) {}; \node[ex] at (5.6,0) {4};
  \node[anchor=west, font=\scriptsize] at (6.2,0) {exones (azul), intrones (gris)};
  \draw[->, thick] (3.0,-0.35) -- (3.0,-0.95) node[midway, right] {transcripción};
  \node[anchor=east] at (-0.2,-1.3) {pre-ARNm};
  \node[ex, fill=omThm!30] at (0.5,-1.3) {1}; \node[intr] at (1.35,-1.3) {}; \node[ex, fill=omThm!30] at (2.2,-1.3) {2}; \node[intr] at (3.05,-1.3) {}; \node[ex, fill=omThm!30] at (3.9,-1.3) {3}; \node[intr] at (4.75,-1.3) {}; \node[ex, fill=omThm!30] at (5.6,-1.3) {4};
  \draw[->, thick] (2.0,-1.65) -- (1.0,-2.35) node[midway, left] {maduración};
  \draw[->, thick] (4.0,-1.65) -- (6.0,-2.35) node[midway, right] {maduración alternativa};
  \node[anchor=east] at (-0.6,-2.7) {ARNm A};
  \node[ex, fill=omThm!30] at (0.0,-2.7) {1}; \node[ex, fill=omThm!30] at (1.0,-2.7) {2}; \node[ex, fill=omThm!30] at (2.0,-2.7) {3}; \node[ex, fill=omThm!30] at (3.0,-2.7) {4};
  \node[anchor=west] at (3.6,-2.7) {proteína A};
  \node[anchor=east] at (5.4,-2.7) {};
  \node[ex, fill=omThm!30] at (6.0,-2.7) {1}; \node[ex, fill=omThm!30] at (7.0,-2.7) {2}; \node[ex, fill=omThm!30] at (8.0,-2.7) {4};
  \node[anchor=west] at (8.6,-2.7) {proteína B};
\end{tikzpicture}

{\small Un \omterm{def:g10:universal-dna:gene}{gen} eucariota se transcribe entero y después se le quitan
los \omterm{prop:g11:gene-expression:splicing}{intrones}. Conservar todos los exones, o dejar uno fuera, da dos
ARNm y dos \omterm{def:g10:chemistry-of-life:families}{proteínas} a partir del mismo \omterm{def:g10:universal-dna:gene}{gen}.}
\end{omfigure}

\begin{remark}[El sentido de la información]\label{rem:g11:gene-expression:flow}
El \omterm{def:g10:universal-dna:information}{ADN} se transcribe en \omterm{def:g11:gene-expression:transcription}{ARN}, el \omterm{def:g11:gene-expression:transcription}{ARN} se traduce en \omterm{def:g11:gene-expression:protein}{proteína} y la
\omterm{def:g11:gene-expression:protein}{proteína} hace el rasgo: la información fluye en un solo sentido. Una
\omterm{def:g11:gene-expression:protein}{proteína} cambiada no reescribe nunca el \omterm{def:g10:universal-dna:gene}{gen}, y por eso toda una vida de
\omterm{def:g10:sport-and-health:training}{entrenamiento} o de quemaduras solares no se hereda, y por eso solo las
\omterm{def:g11:mutations:mutation}{mutaciones} del \omterm{def:g10:universal-dna:information}{ADN} (\cref{ch:g11:mutations}) cambian lo que recibe la
generación siguiente. Los mismos tres pasos funcionan en todas las
\omterm{def:g10:cells-common-unit:cell}{células} de todos los organismos y con el mismo código: la
universalidad del \cref{ch:g10:universal-dna}, ahora hasta la última
palabra.
\end{remark}

\section{Ejercicios}

\begin{exercise}[$\star$]\label{exo:g11:gene-expression:1}
Da tres diferencias entre el \omterm{def:g10:universal-dna:information}{ADN} y el \omterm{def:g11:gene-expression:transcription}{ARN}.
\end{exercise}

\begin{exercise}[$\star$]\label{exo:g11:gene-expression:2}
Transcribe la hebra molde \texttt{3'-TAC CGA AAA ATT-5'} a ARNm.
\end{exercise}

\begin{exercise}[$\star$]\label{exo:g11:gene-expression:3}
Traduce el ARNm \texttt{AUG GGC AAA UGU UAA} con la tabla del código.
\end{exercise}

\begin{exercise}[$\star$]\label{exo:g11:gene-expression:4}
¿Cuáles son los papeles del \omterm{def:g10:cells-common-unit:organelle}{ribosoma} y de los \omterm{def:g11:gene-expression:transcription}{ARN} de transferencia en
la \omterm{prop:g11:gene-expression:translation}{traducción}?
\end{exercise}

\begin{exercise}[$\star$]\label{exo:g11:gene-expression:5}
¿Por qué el código debe usar al menos tres \omterm{def:g10:universal-dna:nucleotide}{nucleótidos} por \omterm{def:g10:chemistry-of-life:families}{aminoácido}?
\end{exercise}

\begin{exercise}[$\star\star$]\label{exo:g11:gene-expression:6}
Una \omterm{def:g11:gene-expression:protein}{proteína} tiene 412 \omterm{def:g10:chemistry-of-life:families}{aminoácidos}. ¿Cuál es la longitud mínima de la
parte codificante de su ARNm, con el \omterm{def:g11:gene-expression:code}{codón} de stop incluido, y la del
\omterm{def:g10:universal-dna:gene}{gen} correspondiente en pares de \omterm{def:g10:universal-dna:nucleotide}{bases}?
\end{exercise}

\begin{exercise}[$\star\star$]\label{exo:g11:gene-expression:7}
La hebra codificante \texttt{ATG AAA GGC TGG TAA} muta a
\texttt{ATG AAA GGC TGA TAA}. Da las dos \omterm{def:g10:chemistry-of-life:families}{proteínas} y clasifica la
\omterm{def:g11:mutations:mutation}{mutación}.
\end{exercise}

\begin{exercise}[$\star\star$]\label{exo:g11:gene-expression:8}
Con la misma secuencia de partida, mutada a \texttt{ATG AAG GGC TGG TAA}
y a \texttt{ATG AAA GGG CTG GTA A}. Clasifica cada una y da las
\omterm{def:g10:chemistry-of-life:families}{proteínas}.
\end{exercise}

\begin{exercise}[$\star\star$]\label{exo:g11:gene-expression:9}
Explica por qué una sustitución en la tercera posición de un \omterm{def:g11:gene-expression:code}{codón} es a
menudo silenciosa, usando la tabla.
\end{exercise}

\begin{exercise}[$\star\star$]\label{exo:g11:gene-expression:10}
El \omterm{def:g11:gene-expression:transcription}{ARN} poli-UC (\texttt{UCUCUCUC}\dots) da una \omterm{def:g11:gene-expression:protein}{proteína} que alterna
serina y leucina. Muestra que eso es coherente con un código de tres
letras y úsalo para asignar dos \omterm{def:g11:gene-expression:code}{codones}.
\end{exercise}

\begin{exercise}[$\star\star$]\label{exo:g11:gene-expression:11}
Un \omterm{def:g10:universal-dna:gene}{gen} humano de \num{30000} pares de \omterm{def:g10:universal-dna:nucleotide}{bases} produce una \omterm{def:g11:gene-expression:protein}{proteína} de 500
\omterm{def:g10:chemistry-of-life:families}{aminoácidos}. Explica la discrepancia.
\end{exercise}

\begin{exercise}[$\star\star\star$]\label{exo:g11:gene-expression:12}
Un fármaco bloquea los \omterm{def:g10:cells-common-unit:organelle}{ribosomas} bacterianos y no los humanos. Explica
por qué puede ser un antibiótico y qué implica su existencia sobre los
\omterm{def:g10:cells-common-unit:organelle}{ribosomas} de ambos grupos.
\end{exercise}

\begin{exercise}[$\star\star\star$]\label{exo:g11:gene-expression:13}
Explica cómo pudo leer un ratón el \omterm{def:g10:universal-dna:gene}{gen} de medusa del
\cref{ch:g10:universal-dna}, con el vocabulario de este capítulo, y qué
ocurriría si el ratón usara otro código.
\end{exercise}

\begin{exercise}[$\star\star\star$]\label{exo:g11:gene-expression:14}
Una \omterm{def:g11:mutations:mutation}{mutación} cambia el ARNt cuyo anticodón lee UAG de modo que lleve un
\omterm{def:g10:chemistry-of-life:families}{aminoácido} en lugar de detener la lectura. Predice su efecto sobre las
\omterm{def:g10:chemistry-of-life:families}{proteínas} de la \omterm{def:g10:cells-common-unit:cell}{célula} en general y sobre un \omterm{def:g10:universal-dna:gene}{gen} mutante que lleve un
UAG prematuro.
\end{exercise}

\begin{exercise}[$\star\star\star$]\label{exo:g11:gene-expression:15}
Discute por qué un desplazamiento del marco de lectura cerca del
principio de un \omterm{def:g10:universal-dna:gene}{gen} es casi siempre peor que uno cerca de su final, y
por qué una \omterm{def:g11:mutations:mutation}{mutación} de cambio de sentido puede ir de inofensiva a
mortal.
\end{exercise}

\section{Problema: Un gen leído de cuatro maneras}

\begin{problem}[{Problema de fin de semana --- un gen corto transcrito,
traducido, mutado tres veces y cronometrado: del ADN a los aminoácidos,
y los sesenta segundos que cuesta una proteína}]
\label{pb:g11:gene-expression:1}
La hebra codificante de un \omterm{def:g10:universal-dna:gene}{gen} corto dice:
\begin{center}
\texttt{ATG GCT TGG AAA CCC GAA TAC GGT CAT TTC TGA}
\end{center}

\medskip
\noindent\textbf{Parte I --- Leer el \omterm{def:g10:universal-dna:gene}{gen}.}
\begin{enumerate}
  \item Escribe la hebra molde.
  \item Escribe el ARNm transcrito a partir del \omterm{def:g10:universal-dna:gene}{gen}.
  \item Tradúcelo \omterm{def:g11:gene-expression:code}{codón} a \omterm{def:g11:gene-expression:code}{codón} y da la secuencia de \omterm{def:g10:chemistry-of-life:families}{aminoácidos} de la
        \omterm{def:g11:gene-expression:protein}{proteína}.
  \item ¿Cuántos \omterm{def:g10:chemistry-of-life:families}{aminoácidos} tiene la \omterm{def:g11:gene-expression:protein}{proteína}, y cuántos \omterm{def:g10:universal-dna:nucleotide}{nucleótidos}
        del ARNm se usan para especificarlos, con el stop incluido?
  \item ¿Qué \omterm{def:g11:gene-expression:code}{codón} inició la lectura, y qué fija el marco de lectura de
        todos los demás?
\end{enumerate}

\medskip
\noindent\textbf{Parte II --- Tres mutantes.}
\begin{enumerate}[resume]
  \item Mutante 1: \texttt{ATG GCT TGG AAA CCG GAA TAC GGT CAT TTC TGA}.
        Da la \omterm{def:g11:gene-expression:protein}{proteína} y clasifica la \omterm{def:g11:mutations:mutation}{mutación}.
  \item Mutante 2: \texttt{ATG GCT TGG AAA CCC GAA TAC GGT CAT TTC TGA}
        con el noveno \omterm{def:g11:gene-expression:code}{codón} \texttt{CAT} sustituido por \texttt{CGT}.
        Da la \omterm{def:g11:gene-expression:protein}{proteína} y clasifícala.
  \item Mutante 3: \texttt{ATG GCT TGA AAA CCC GAA TAC GGT CAT TTC TGA}.
        Da la \omterm{def:g11:gene-expression:protein}{proteína} y clasifícala. ¿Qué única sustitución la
        produjo?
  \item Mutante 4: se inserta una T tras el sexto \omterm{def:g10:universal-dna:nucleotide}{nucleótido}:
        \texttt{ATG GCT TTG GAA ACC CGA ATA CGG TCA TTT CTG A}. Da la
        \omterm{def:g11:gene-expression:protein}{proteína} y clasifícala.
  \item Ordena los cuatro mutantes de menos a más perturbador para la
        \omterm{def:g11:gene-expression:protein}{proteína} y justifícalo.
\end{enumerate}

\medskip
\noindent\textbf{Parte III --- Por qué el código es como es.}
\begin{enumerate}[resume]
  \item ¿Cuántos \omterm{def:g11:gene-expression:code}{codones} especifican la leucina? ¿Y la serina? ¿Y el
        triptófano? ¿Qué \omterm{def:g10:chemistry-of-life:families}{aminoácido} tiene más probabilidades de cambiar
        por una sustitución al azar en su \omterm{def:g11:gene-expression:code}{codón}, y cuál menos?
  \item Calcula la fracción de todas las sustituciones en la tercera
        posición de un \omterm{def:g11:gene-expression:code}{codón} de leucina \texttt{CUx} que resultan
        silenciosas.
  \item Explica por qué un código de dos letras no podría funcionar y
        por qué no hace falta uno de cuatro.
  \item El poli-U da polifenilalanina y el poli-A polilisina. ¿Qué
        entradas de la tabla establecen esos dos experimentos?
  \item El código es el mismo en la bacteria y en un humano. Enuncia la
        consecuencia para la \omterm{prop:g10:universal-dna:universal}{transgénesis} y la consecuencia para el
        parentesco.
\end{enumerate}

\medskip
\noindent\textbf{Parte IV --- Cronometrar la fábrica.}
Un \omterm{def:g10:cells-common-unit:organelle}{ribosoma} añade 10 \omterm{def:g10:chemistry-of-life:families}{aminoácidos} por segundo y los \omterm{def:g10:cells-common-unit:organelle}{ribosomas} se siguen
unos a otros a 80 \omterm{def:g10:universal-dna:nucleotide}{nucleótidos} de distancia a lo largo de un ARNm; cada
ARNm vive 2 horas.
\begin{enumerate}[resume]
  \item ¿Cuánto tarda un \omterm{def:g10:cells-common-unit:organelle}{ribosoma} en traducir la \omterm{def:g11:gene-expression:protein}{proteína} de este \omterm{def:g10:universal-dna:gene}{gen}?
        ¿Y una \omterm{def:g11:gene-expression:protein}{proteína} de 600 \omterm{def:g10:chemistry-of-life:families}{aminoácidos}?
  \item ¿Cuántos \omterm{def:g10:cells-common-unit:organelle}{ribosomas} pueden leer a la vez un ARNm de 1800
        \omterm{def:g10:universal-dna:nucleotide}{nucleótidos}?
  \item Si arranca un \omterm{def:g10:cells-common-unit:organelle}{ribosoma} nuevo cada 8 segundos, ¿cuántas copias
        de la \omterm{def:g11:gene-expression:protein}{proteína} rinde un ARNm a lo largo de su vida?
  \item Una bacteria necesita \num{2000} copias de una \omterm{def:g10:cell-metabolism:metabolism}{enzima} en 10
        minutos. ¿Cuántos ARNm del \omterm{def:g10:universal-dna:gene}{gen}, transcritos a la vez, hacen
        falta?
  \item Enuncia el resultado: la \omterm{def:g11:gene-expression:protein}{proteína} que codifica el \omterm{def:g10:universal-dna:gene}{gen}, la
        \omterm{def:g11:mutations:mutation}{mutación} de las cuatro que la abolió y el tiempo que necesita
        un \omterm{def:g10:cells-common-unit:organelle}{ribosoma} para construirla.
\end{enumerate}
\end{problem}
